\section{The complete signed support and the absolute twistor charts}


\subsection{Human source and the exact categories}

The source is Alain Connes and Caterina Consani,
\emph{The Absolute Twistor Line and the Geometry of
$\overline{\Spec\Z}$},
\href{https://arxiv.org/abs/2609.00299v1}{arXiv:2609.00299v1}.
The original author file \path{CC.tex} was read at lines
310--413 and 420--1207, including the entire relevant sections on
the signed base, the Lydakis--Day product and proposed coproduct,
the two charts, the geometric involution, complex conjugation,
and odd Frobenius. Exact source locators are
\texttt{F2}, \texttt{sect2}, \texttt{sect2.2}, \texttt{conv},
\texttt{opens}, \texttt{eq:restriction\_maps},
\texttt{def:alpha\_symmetry}, \texttt{thm:twistor\_space\_scheme},
and \texttt{prop:twistor\_frobenius}.

The calculations below first take place in commutative pointed
monoids. Their zero is absorbing and their identity is denoted $1$.
Put
\[
 B=\{0,1,\epsilon\},\qquad\epsilon^2=1,\quad\epsilon\ne1.
 \tag{AT1}\label{at:eq:signed-base}
\]
A signed pointed monoid is supplied with a morphism from $B$.
The target may specialize the sign, but all universal calculations
retain $\epsilon\ne1$. The source identifies spherical monoids'
tensor product with the relative smash product. We therefore write
$\otimes_B$ for that pointed-monoid pushout: in particular
$\epsilon\otimes1=1\otimes\epsilon$. No additive relation is
introduced until an explicitly named ring or semiring receiver.

\subsection{The actual full split base and its two spectra}

Let $L$ be a finite nontrivial bounded distributive lattice. Retain
\[
 G_L(\Z)=\{z_\lambda=(0,\lambda):\lambda\in L\}
       \cup\{\widehat a=(a,1_L):a\in\Z\}.
 \tag{AT2}\label{at:eq:full-semiring}
\]
Its overlap is $e=z_{1_L}=\widehat0$, while
$\tau=z_{0_L}\ne e$. Addition is coefficient addition and lattice
join; multiplication is coefficient multiplication and lattice
meet. Thus $\widehat a z_\lambda=z_\lambda$ and
$z_\lambda z_\mu=z_{\lambda\wedge\mu}$. Let $M_L$ be its
underlying multiplicative pointed monoid with zero $\tau$.
There is an injective signed-base map
\[
 B\longrightarrow M_L,
 \qquad 0\mapsto\tau,\quad1\mapsto\widehat1,
 \quad\epsilon\mapsto\widehat{-1}.
 \tag{AT3}\label{at:eq:base-map}
\]
In particular $\epsilon z_\lambda=z_\lambda$ for every support
label. This relation is retained below, including for $e$.

For precision a prime of a pointed monoid is a proper ideal whose
complement is multiplicatively closed; additive closure is not
part of this definition. The entire pointed-monoid spectrum of
$M_L$ can be computed:
\[
 \begin{aligned}
 P_c&=\{z_\lambda:c\not\le\lambda\},&&0_L<c\in L,\\
 Q_E&=\{z_\lambda:\lambda\in L\}
  \cup\{\widehat a:a\ne0,\ p\mid a\text{ for some }p\in E\},
  &&E\subset\{\text{rational primes}\}.
 \end{aligned}
 \tag{AT4}\label{at:eq:monoid-primes}
\]
To prove exhaustion, an ideal containing any supported coefficient
contains $e$ by multiplication by $e$, and hence every support
label. Its remaining prime condition is exactly that on the free
commutative monoid of positive prime factors, with signs units;
unique integer factorization gives $Q_E$. An ideal omitting $e$
contains no supported coefficient. Its support labels form a
downset whose complement is a nonempty meet-closed up-set. Because
$L$ is finite, that complement is $\mathord\uparrow c$ for the
meet $c$ of its elements, and $c\ne0_L$. This gives $P_c$.
Both displayed families directly satisfy the ideal and prime
conditions, proving the classification.

The original semiring spectrum embeds topologically in this monoid
spectrum by retaining exactly the additively closed primes. In
the support family these are the $P_c$ with $c$ join prime, or
equivalently join irreducible in a finite distributive lattice.
Indeed join closure of its complement's complement is exactly
$c\le\lambda\vee\mu\Rightarrow c\le\lambda$ or $c\le\mu$.
In the arithmetic family they are $Q_\varnothing$ and
$Q_{\{p\}}$: two different primes in $E$ generate $1$ additively
by their integer Bezout identity, and a singleton gives the usual
prime ideal. Its subspace topology is the original one, since
the basic opens in both cases are the same tests $a\notin P$.
Thus the full semiring spectrum
$\SpecLat L\triangleright\Spec\Z$ has been preserved by an
explicit embedding, not identified with the larger monoid spectrum.

\subsection{Adjoining the imaginary generator without losing a support label}

For any signed pointed monoid $M$, set
\[
 A_M=M[J]/(J^2=\epsilon).
 \tag{AT5}\label{at:eq:imaginary-extension}
\]
There is an explicit normal form: its elements are $(m,b)$ with
$m\in M$ and $b\in\{0,1\}$, identifying only $(0,0)$ and
$(0,1)$. Write $(m,b)=mJ^b$. Multiplication is
\[
 (m,b)(n,c)=(mn\epsilon^{bc},\ b+c\bmod2).
 \tag{AT6}\label{at:eq:normal-product}
\]
The carry identity modulo two proves associativity, and the
formula has $J=(1,1)$ with $J^2=\epsilon$. A map out of this
monoid is uniquely determined by a signed map out of $M$ and
a root of $j^2=\epsilon$; reduction of every power of $j$ proves
the universal property. This proves the normal form, including
the injectivity of $M\to A_M$.

For $M=M_L$ it follows that above every nonbottom support label
there are two distinct elements
\[
 z_\lambda,\quad z_\lambda J\qquad(\lambda\ne0_L),
 \quad (z_\lambda J)^2=z_\lambda.
 \tag{AT7}\label{at:eq:support-tags}
\]
The equality uses $\epsilon z_\lambda=z_\lambda$; it does not
identify these two elements. There is a canonical map to an
actual Gaussian coefficient receiver,
\[
 A_{M_L}\longrightarrow U(G_L(\Z[i])),\qquad
 \widehat aJ^b\longmapsto\widehat{ai^b},\qquad
 z_\lambda J^b\longmapsto z_\lambda,
 \tag{AT8}\label{at:eq:gaussian-map}
\]
where $U$ denotes forgetting addition. Its image is exactly the
axial submonoid consisting of the support labels and nonzero real
or purely imaginary integer coefficients. Its kernel congruence
is exactly the identifications $z_\lambda J=z_\lambda$ for
$\lambda\ne0_L$: distinct nonzero axial coefficients remain
distinct, and the normal form proves there are no other
identifications. The additive closure of this image is the full
$G_L(\Z[i])$, since $a+bi$ is a sum of its two displayed axes.
This describes both the retained data and the precise additional
quotient in passing to Gaussian coefficients.

There is no signed map $A_B\to M_L$: a root of $j^2=\epsilon$
would be a unit, but the units of $M_L$ are $\widehat1$ and
$\widehat{-1}$, whose squares are $1$. Thus
\eqref{at:eq:imaginary-extension} and \eqref{at:eq:gaussian-map}
are the required root-adjoining maps, rather than an assumed
root in the original integer coefficient monoid.

Adjoining this unit does not change the monoid prime topology.
For every prime of $A_M$, membership of $mJ^b$ is equivalent to
membership of $m$, because $J$ is a unit. Conversely the same
rule extends every prime of $M$ uniquely to one of $A_M$;
the normal form and the fact that $\epsilon$ is a unit prove
well-definedness. Basic opens satisfy $D(mJ^b)=D(m)$. Hence
\[
 \Spec A_M\xrightarrow{\sim}\Spec M
 \tag{AT9}\label{at:eq:spectrum-root}
\]
is a homeomorphism. The extra roots and the extra support tags
therefore have an exact point-functor meaning even where they
do not separate additional monoid primes.

\subsection{The diagonal obstruction and the actual torsor coaction}

For a free signed spatial monoid $B[N]$, its spatial diagonal
$n\mapsto n\otimes n$ together with $\epsilon\mapsto\epsilon$
is defined over the fixed base. The relation $J^2=\epsilon$ is
not a relation internal to an unsigned spatial monoid independent
of that base. The proposed diagonal must therefore be checked
after imposing this relation.

\begin{theorem}[Failure of the binary diagonal on the actual relation]
\label{at:thm:diagonal-failure}
For $A=A_B=B[J]/(J^2=\epsilon)$, no signed morphism
$A\to A\otimes_BA$ can send $J$ to $J_1J_2$.
Indeed
\[
 (J_1J_2)^2=\epsilon^2=1\ne\epsilon.
 \tag{AT10}\label{at:eq:diagonal-defect}
\]
The same obstruction occurs in $A_M$ whenever the image of
$\epsilon$ remains distinct from $1$ in $A_M\otimes_MA_M$.
The raw diagonal is defined after additionally imposing
$\epsilon=1$, but that is a different signed-base receiver.
\end{theorem}
\begin{proof}
A signed morphism must send $J^2$ to the same base element
$\epsilon$. Its proposed value has the square in
\eqref{at:eq:diagonal-defect}. To verify that these elements remain
different in the displayed target, map both $J_1,J_2$ to $i$
in the multiplicative pointed monoid of $\C$, and send
$\epsilon$ to $-1$. This respects the defining tensor relations
and distinguishes $1$ from $\epsilon$. Over $M_L$ the original
coefficient map to $\C$ also distinguishes them, so the same
counterexample applies to the full base extension. After the
specified sign collapse, the equality required for descent
does hold, proving the last assertion.
\end{proof}

The counit obstruction is equally exact: $A_B$ has no $B$-point
because neither $1$ nor $\epsilon$ has square $\epsilon$.
Thus this imaginary-root object cannot have an identity section
for a group law over $B$. This is a defect of applying the free
spatial diagonal to this particular signed relation; it does not
remove the valid coproduct on the free spatial core $B[T^{\Z}]$.

There is instead the genuine group object
\[
 H=B[U]/(U^2=1),\qquad
 \Delta_H(U)=U_1U_2,\quad e_H(U)=1,\quad\iota_H(U)=U^{-1}=U.
 \tag{AT11}\label{at:eq:mu-two}
\]
All base elements are fixed. Direct substitution proves the
group identities. The actual imaginary-root coaction is
\[
 \rho:A_B\longrightarrow A_B\otimes_BH,
 \qquad \rho(J)=J\otimes U.
 \tag{AT12}\label{at:eq:coaction}
\]
Its square is $\epsilon\otimes1$, and applying either composite
coaction gives $J\otimes U\otimes U$, proving coassociativity;
the counit gives $J$ back. Base change gives the same formulas
over every original $M_L$.

\begin{theorem}[The precise torsor identity]
\label{at:thm:torsor}
The morphism of root functors
$(j,u)\mapsto(j,ju)$ is represented by an isomorphism
\[
 A_B\otimes_BA_B\xrightarrow{\sim}A_B\otimes_BH,
 \qquad J_1\mapsto J,\quad J_2\mapsto JU,
 \tag{AT13}\label{at:eq:torsor-identity}
\]
whose inverse sends $J\mapsto J_1$ and
$U\mapsto J_1^{-1}J_2$. For every signed target the roots of
$j^2=\epsilon$, when nonempty, form a free transitive set under
the units $u$ with $u^2=1$.
\end{theorem}
\begin{proof}
All imaginary roots are units because $j^4=1$ and
$j^{-1}=\epsilon j$. Therefore the inverse formula is defined,
and its proposed $U$ has square $\epsilon^{-1}\epsilon=1$.
Composing the two displayed formulas gives the identity on each
generator, proving the isomorphism. For two roots $j_1,j_2$ the
unique translating unit is $j_1^{-1}j_2$, with square one.
This proves the assertion on every target and its base changes.
\end{proof}

Here ``torsor'' means this explicit action and isomorphism;
no unspecified topology on pointed monoids is needed. After
passing to any ring with $\epsilon=-1$, the object is
$\operatorname{Spec}R[J]/(J^2+1)$, finite free of rank two and
faithfully flat over $R$, with precisely the same $\mu_2$ torsor
identity. Faithful flatness follows because its module is $R^2$.
Thus it is also a torsor in the faithfully flat ring-scheme sense.

The canonical ternary operation is the torsor heap
\[
 h(j_1,j_2,j_3)=j_1j_2^{-1}j_3,
 \qquad h^*(J)=J_1J_2^{-1}J_3.
 \tag{AT14}\label{at:eq:heap}
\]
Its square is $\epsilon\epsilon^{-1}\epsilon=\epsilon$.
It satisfies $h(x,y,y)=x=h(y,y,x)$ and
$h(h(x,y,z),u,v)=h(x,y,h(z,u,v))$ by multiplication of units.
Every odd raw product of roots also has square $\epsilon$;
every even raw product has square one. These operations are
different: $j_1j_2j_3=\epsilon h(j_1,j_2,j_3)$. Its odd
diagonal is the actual power $j^{2k+1}=\epsilon^kj$.
Choosing a root $j_0$ after a specified base extension gives
the binary law $j_1j_0^{-1}j_2$ with identity $j_0$, thereby
identifying the torsor with $\mu_2$ after that choice. No such
choice was present over $B$.

\subsection{The actual charts, overlap, and geometric involution}

Use exactly the source charts, now in pointed-monoid notation:
\[
 A_+=A_B[T],\qquad A_-=A_B[S],\qquad
 A_\circ=A_B[T,T^{-1}],\qquad
 S=T^{-1},\quad J_+=J,\quad J_-=\epsilon J.
 \tag{AT15}\label{at:eq:charts}
\]
The coaction \eqref{at:eq:coaction} extends by fixing the spatial
coordinate. It respects the overlap because
$\rho(J_-)=\epsilon JU=\rho(\epsilon J)$.
The source's geometric involution is
\[
 \begin{gathered}\alpha^*(T)=\epsilon T^{-1},\qquad
 \alpha^*(J)=\epsilon J=J^{-1};\\
 \quad
 \alpha_+^*(S)=\epsilon T,\quad\alpha_+^*(J_-)=J_+,\\
 \quad
 \alpha_-^*(T)=\epsilon S,\quad\alpha_-^*(J_+)=J_-.
 \end{gathered}\tag{AT16}\label{at:eq:alpha}
\]
It squares to the identity on every generator. For instance the
overlap of $\alpha_+^*(J_-)$ is $J$, while
$\alpha^*(\epsilon J)=\epsilon^2J=J$. This verifies the
nontrivial sign in the chart compatibility. The torsor coaction
commutes with $\alpha$: on $J$ both give $\epsilon JU$, and
on $T$ both give $\epsilon T^{-1}$. The same formulas hold
after the full base change from $B$ to $M_L$, with $\alpha$
fixing every element of $M_L$.

For the pure spherical generic chart, its monomials have the
form $mJ^bT^r$, $m\in\{1,\epsilon\}$, $b=0,1$, $r\in\Z$.
Invariance under \eqref{at:eq:alpha} forces $r=0$ and then $b=0$.
Thus its invariant pointed monoid is exactly $B$, as used in
the source's gluing argument. After the actual full base change
its invariant pointed monoid is instead
\[
 (A_{M_L}[T,T^{-1}])^\alpha
    =M_L\ \cup\ \{z_\lambda J:0_L<\lambda\in L\}.
 \tag{AT17}\label{at:eq:full-invariants}
\]
Indeed the exponent comparison still forces $r=0$, and an odd
$J$ exponent is invariant precisely when $\epsilon m=m$.
By \eqref{at:eq:full-semiring}, these nonzero $m$ are exactly
the support labels. This proves the exact changed invariant
receiver while retaining the source's original pure-base result.

For completeness the chart spectra also keep the entire original
base. Every prime of $M[T]$ is determined by its contraction $P$
and whether it contains $T$. If it omits $T$, then
$mT^r\in P'$ precisely when $m\in P$; if it contains $T$,
all positive exponents belong and the exponent-zero part is $P$.
These rules give primes and prove exhaustion. Basic opens show
the resulting spectrum is the product of $\Spec M$ with the
two-point affine monoid line. Gluing the two localizations in
\eqref{at:eq:charts}, and using \eqref{at:eq:spectrum-root}, gives
the product with the three-point projective monoid line.
The injection of the full semiring spectrum from
\eqref{at:eq:monoid-primes} is retained in this construction.
This is a base change of the actual charts, not an identification
with the source's separate global amalgam.

\subsection{Odd powers, fixed points, and the exact Frobenius bridge}

For $n\ge1$, power $x\mapsto x^n$ is a pointed-monoid
endomorphism of any commutative pointed monoid. It fixes the
base \eqref{at:eq:signed-base} precisely when $n$ is odd. On the
source charts the resulting maps are
\[
 F_n^*(T)=T^n,\qquad F_n^*(J)=J^n=\epsilon^{(n-1)/2}J,
 \qquad n\text{ odd}.
 \tag{AT18}\label{at:eq:odd-powers}
\]
They respect $J^2=\epsilon$, every overlap, and $\alpha$:
for example $\alpha^*(T^n)=\epsilon^nT^{-n}
=\epsilon T^{-n}$, and $F_n^*(\epsilon J)=\epsilon J^n$.
They compose as $F_mF_n=F_{mn}$. On the torsor the equivariance
identity is $\rho F_n=(F_n\otimes[n])\rho$, with
$[n](U)=U^n=U$ for odd $n$. Thus these powers are defined without
the invalid binary imaginary diagonal.

After adjoining the full original integer coefficient monoid there
are two different exact maps. The relative map fixes $M_L$ and
uses \eqref{at:eq:odd-powers} on $T,J$. The absolute pointed-monoid
power also sends each $m\in M_L$ to $m^n$:
\[
 F_{n,\mathrm{rel}}(mJ^bT^r)=mJ^{bn}T^{rn},\qquad
 F_{n,\mathrm{abs}}(mJ^bT^r)=m^nJ^{bn}T^{rn}.
 \tag{AT19}\label{at:eq:two-powers}
\]
Both are signed for odd $n$. In the Gaussian polynomial receiver
of \eqref{at:eq:gaussian-map}, the relative map extends to the actual
ring homomorphism fixing $\Z$ and sending $i\mapsto i^n$,
$T\mapsto T^n$, and hence to the full split semiring while fixing
every support label. The absolute map is the multiplicative power
on that receiver. In characteristic zero it is not additive for
$n>1$, since $(1+1)^n=2^n\ne2=1^n+1^n$. This is the exact
comparison between the two coefficient operations.

Power induces the identity on every monoid prime spectrum:
$x^n\in P$ if and only if $x\in P$, by primality and induction.
For the relative chart map, the preceding chart-prime classification
gives the same identity: it fixes the base prime and the condition
whether $T$ belongs. This does not imply triviality on field-valued
points; those are calculated below. In particular the original
semiring prime subspace of $G_L(\Z)$ is fixed by absolute powers
as a topological space, even though those maps are not semiring
endomorphisms in characteristic zero.

The distinction is visible in the actual Gaussian ring spectrum.
For the relative $n=3$ substitution on $\Z[i][T]$, the inverse
image of the prime $(T-2)$ is $(T-8)$: its quotient test is
$f(i,T)\mapsto f(-i,8)$, and Gaussian conjugation is an
automorphism. For absolute cubing, inverse image of the same
prime as a multiplicative ideal is $(T-2)$ itself. Thus the
identity statement about pointed-monoid powers has not been
extended incorrectly to the additive substitution receiver.

The joint fixed set of all odd powers on the pure source chart is
$B$: already the third power forces the spatial exponent to be
zero and removes the odd imaginary exponent. On the full
absolute base change, the exact fixed set is
\[
 \{\tau,\widehat1,\widehat{-1}\}
 \ \cup\ \{z_\lambda J^b:0_L<\lambda\in L,\ b=0,1\}.
 \tag{AT20}\label{at:eq:full-power-fixed}
\]
The normal form proves this directly. A nonzero integer coefficient
fixed by the third power is $1$ or $-1$. An odd imaginary exponent
on such a coefficient would require $-a^3=a$, impossible for a
nonzero integer. By contrast every displayed support element is
fixed by all odd powers because $z_\lambda^n=z_\lambda$ and
$\epsilon z_\lambda=z_\lambda$. This includes both tags over
the original supported zero $e$.
For relative odd powers the integer coefficients are fixed as
well, and the joint fixed set is exactly the right-hand side of
\eqref{at:eq:full-invariants}. This gives both versions of the
fixed-point comparison.

\subsection{Complex points, the twistor quotient, and the action that descends}

In a field-valued point the sign is evaluated as $-1$. If the field
has no root of $-1$, the imaginary charts have no points. If it
has characteristic different from two and roots $\pm i$, the
chart overlap is exactly
\[
 (t,j)_+\sim(t^{-1},-j)_-.
 \tag{AT21}\label{at:eq:point-overlap}
\]
Thus there are two projective lines, with the first joining the
$(+,i)$ and $(-,-i)$ charts. In characteristic two the unique
field-valued root is $1$ and the point sets form a single
projective line. This is a statement about field-valued points,
not a claim that the characteristic-two root scheme is reduced.

On the complex generic open the original involution is
\[
 \alpha(z,j)=(-z^{-1},-j),\qquad
 c(z,j)=(\overline z,-j).
 \tag{AT22}\label{at:eq:point-involutions}
\]
It interchanges the two spheres, and each orbit has a unique
representative on the first sphere. Taking that representative
after conjugation gives the exact real structure
$\sigma(z)=-1/\overline z$, including $0\leftrightarrow\infty$.
It has no fixed point, since a nonzero finite fixed point would
satisfy $|z|^2=-1$ and the endpoints are exchanged.

There is a necessary descent check for the spatial action. Scaling
both root sheets by $a_\lambda(z,j)=(\lambda z,j)$ satisfies
$\alpha a_\lambda=a_{\lambda^{-1}}\alpha$; it does not
preserve the $\alpha$-orbits in general. Indeed the representatives
$(z,i)$ and $(-z^{-1},-i)$ scale to
$(\lambda z,i)$ and $(-\lambda z^{-1},-i)$, which lie in one
$\alpha$-orbit exactly when $\lambda=\lambda^{-1}$.
The exact action which does descend is
\[
 \widetilde a_\lambda(z,i)=(\lambda z,i),\qquad
 \widetilde a_\lambda(z,-i)=(\lambda^{-1}z,-i).
 \tag{AT23}\label{at:eq:opposite-scaling}
\]
It is an action, extends across both poles, and commutes with
$\alpha$ by direct substitution in \eqref{at:eq:point-involutions}.
On the first representative sphere it gives the usual action
$z\mapsto\lambda z$. Complex conjugation transforms its parameter
by $\lambda\mapsto\overline\lambda^{-1}$, so
\[
 \sigma(\lambda z)=\overline\lambda^{-1}\sigma(z).
 \tag{AT24}\label{at:eq:twisted-equivariance}
\]
This is the actual twistor equivariance, now with its descent map.

There is a real algebraic group recording exactly that parameter
descent: the norm-one torus
\[
 \mathcal T=\Spec\R[a,b]/(a^2+b^2-1),
 \quad(a,b)(c,d)=(ac-bd,ad+bc).
 \tag{AT25}\label{at:eq:real-torus}
\]
Over $\C$, the coordinate $\lambda=a+ib$ identifies it with
$\G_m$, with inverse $a=(\lambda+\lambda^{-1})/2$,
$b=(\lambda-\lambda^{-1})/(2i)$. The real structure on this
complex parameter is $\lambda\mapsto\overline\lambda^{-1}$,
as those formulas show. This realizes \eqref{at:eq:opposite-scaling}
and \eqref{at:eq:twisted-equivariance} after the specified real
additive extension. Division by two is stated explicitly; no
identification of $\mu_2$ with a constant two-point group in
characteristic two has been used.

For an odd $n$, let $\chi_4(n)=(-1)^{(n-1)/2}$. On the two
sheets the actual map is $(z,j)\mapsto(z^n,j^n)$. Taking the
first-sheet representative gives precisely
\[
 f_n(z)=\begin{cases}z^n,&\chi_4(n)=1,\\
                      -z^{-n},&\chi_4(n)=-1.
        \end{cases}
 \tag{AT26}\label{at:eq:twistor-power}
\]
Both formulas extend over zero and infinity. Put $a(z)=-1/z$.
For odd $n$ it commutes with $z\mapsto z^n$, and $a^2=1$.
Since $\chi_4(mn)=\chi_4(m)\chi_4(n)$, these maps satisfy
$f_mf_n=f_{mn}$. They commute with $\sigma$ by conjugating their
real coefficients and using the same odd sign. Their precise
equivariance with the descended torus action is
\[
 f_n(\lambda z)=\lambda^{n\chi_4(n)}f_n(z).
 \tag{AT27}\label{at:eq:power-character}
\]
Thus the second case reverses the torus character and exchanges
the poles. Its map of projective curves still has degree $n$:
the pullback of the zero divisor is $n[0]$ in the first case
and $n[\infty]$ in the second, with the opposite formula for the
infinity divisor. These explicit divisor and character maps are
the proved comparisons; no assertion about an unspecified Hodge
bundle is needed.

\subsection{The additive invariant conic and the prime two}

There is also an exact additive receiver for the source's generic
involution. In the ring $R_\circ=\Z[i][T,T^{-1}]$ put
$\alpha(i)=-i$ and $\alpha(T)=-T^{-1}$. Define
\[
 u=T-T^{-1},\qquad v=i(T+T^{-1}).
 \tag{AT28}\label{at:eq:conic-coordinates}
\]
Both are invariant and $u^2+v^2=-4$.

\begin{theorem}[The entire additive invariant receiver]
\label{at:thm:invariant-conic}
There are exact identifications
\[
 R_\circ^\alpha=\Z[u,v]/(u^2+v^2+4),\qquad
 G_L(R_\circ)^\alpha=G_L\bigl(\Z[u,v]/(u^2+v^2+4)\bigr).
 \tag{AT29}\label{at:eq:invariant-conic}
\]
They retain every lattice label. The monoid invariant map from
\eqref{at:eq:full-invariants} lands in the original integer
coefficient and support part; the new $u,v$ are actual sums
in this additive receiver.
\end{theorem}
\begin{proof}
An invariant Laurent polynomial has integer constant coefficient
and, for every $r>0$, a pair
$a_rT^r+(-1)^r\overline{a_r}T^{-r}$, with $a_r\in\Z[i]$.
Its real and imaginary integer parts are respectively multiples of
\[
 F_r=T^r+(-1)^rT^{-r},\qquad
 G_r=i(T^r-(-1)^rT^{-r}).
\]
Both satisfy $X_{r+1}=uX_r+X_{r-1}$, with
$F_0=2,F_1=u,G_0=0,G_1=v$. Thus all invariants are generated
by $u,v$ and integer constants. The indicated relation holds.
It is the entire kernel: because it is monic in $v^2$, any
polynomial reduces to $A(u)+vB(u)$. In the Laurent receiver the
first summand has integer coefficients and the second purely
imaginary integer coefficients, so their sum vanishes only if
both vanish. The element $u=T-T^{-1}$ is transcendental over
$\Z$ (its powers have distinct highest $T$ exponents), and
$T+T^{-1}$ is nonzero in this domain. Hence $A=B=0$.
This proves the first identification. The involution fixes every
$z_\lambda$ and acts on a supported coefficient exactly by its
ring involution, proving the second. Finally the normal-form
fixed elements in \eqref{at:eq:full-invariants} map either to their
integer coefficient or to their original support label under
\eqref{at:eq:gaussian-map}, proving the stated comparison.
\end{proof}

The conic in \eqref{at:eq:invariant-conic} has no real points:
real squares cannot sum to $-4$. It is an explicitly computed
additive enlargement of the invariant receiver, not a claim
that invariants commute with adding ring operations. Its proof
also preserves the exact integer constant $4$.

Over $\C$ the generic conic is the actual folded generic curve,
with inverse coordinate
\[
 z=\frac{u-iv}{2},\qquad z^{-1}=-\frac{u+iv}{2}.
 \tag{AT29a}\label{at:eq:conic-inverse}
\]
Their product is one by the conic equation. For each choice
$j\in\{i,-i\}$ the corresponding unfolded coordinate is
$T=(u-jv)/2$; the other choice gives $-T^{-1}$, so these are
exactly one $\alpha$-orbit. This proves the quotient statement
directly on all generic complex points. Conjugating the real
coordinates $u,v$ induces $z\mapsto-1/\overline z$, as
substitution in \eqref{at:eq:conic-inverse} verifies.
The torus \eqref{at:eq:real-torus} acts on this conic by
\[
 (u,v)\longmapsto(au+bv,\ av-bu),\qquad a^2+b^2=1.
 \tag{AT29b}\label{at:eq:conic-action}
\]
It preserves $u^2+v^2+4$ and, in \eqref{at:eq:conic-inverse},
becomes $z\mapsto(a+ib)z$. Thus the repaired scalar descent
has a concrete additive-coordinate action with the same signs.
For odd $n$, the relative Frobenius acts on these invariant
coordinates by $u\mapsto F_n$ and
$v\mapsto\chi_4(n)G_n$, where the polynomials $F_n,G_n$
were constructed by the exact recurrence in the proof above.
This follows by substituting $T\mapsto T^n$ and $i\mapsto i^n$,
and is the conic expression of \eqref{at:eq:twistor-power}.

Finally the root torsor's reduction at two is an actual scheme
calculation. In $\Z[i]$,
$(1+i)^2=2i$, so $(2)=(1+i)^2$ as ideals, and
\[
 \Z[i]/2\Z[i]\simeq\mathbb F_2[J]/((J+1)^2).
 \tag{AT30}\label{at:eq:two-fibre}
\]
The ring has a nonzero nilpotent $J+1$ although it has only one
field-valued root. At an odd prime the polynomial $J^2+1$ has
no common root with its derivative $2J$, so its fibre is reduced
and the degree-two algebra is etale there. The trace matrix in
the basis $1,i$ is $\operatorname{diag}(2,-2)$, with determinant
$-4$, in agreement with these exact fibres. This supplies the
arithmetical relation to the prime two. The restriction of
pointed-monoid powers to odd integers was separately proved by
fixing $\epsilon$; neither calculation is substituted for the
other.

\begin{figure}[ht]
\centering
\begin{tikzcd}[column sep=large,row sep=large]
 M_L\arrow[r,hook]\arrow[d,hook]
 &A_{M_L}\arrow[d]\arrow[r,"\rho"]
 &A_{M_L}\otimes_B H\\
 U(G_L(\Z[i]))\arrow[r,equal]
 &U(G_L(\Z[i]))
\end{tikzcd}
\caption{The actual signed bridge retains the full integer split
base on the left. The upper inclusion adjoins $J^2=\epsilon$.
The middle downward map is \eqref{at:eq:gaussian-map}: it identifies
only the two imaginary tags above each nonbottom support label
within its axial image. Its right-hand coaction is the valid
torsor map \eqref{at:eq:coaction}, not the binary diagonal ruled
out in Theorem~\ref{at:thm:diagonal-failure}. The full Gaussian
receiver is obtained by taking the stated additive closure.}
\end{figure}

The free spatial coproduct, imaginary torsor, ternary operation,
odd powers, full-support fixed points, Gaussian coefficient maps,
and twistor quotient have all been related by their explicit
morphisms. The conclusions use the author's signed chart relations
and field hypotheses exactly. No Riemann hypothesis, smoothness
of a global arithmetic curve, or spectral purity conclusion is
asserted from these constructions.

